%------------------------------------------------------------------------------%
\begin{question}
  Sabendo que os autovalores da matriz
  \[
    B =
    \begin{pmatrix}
      3 & 1 & 1 \\
      1 & 3 & 1 \\
      1 & 1 & 3
    \end{pmatrix}
  \]
  são
  \(
    \lambda_{1}=5
  \)
  e
  \(
    \lambda_{2}=2
  \)

  \bigskip
  Determine as matrizes $P$ e $D$ tais que
  \(
    B = P D P^t
  \)
\end{question}

%------------------------------------------------------------------------------%
\begin{resolution}{Autovalores}
  Calculamos o polinômio característico
  \begin{align*}
      p(\lambda)
      &=\det(B-\lambda I)\\
      &=
      \begin{vmatrix}
          3-\lambda & 1 & 1\\
          1 & 3-\lambda & 1\\
          1 & 1 & 3-\lambda
      \end{vmatrix}\\
      &=(5-\lambda)(2-\lambda)^2
  \end{align*}

  Portanto
  \[
      \lambda_{1}=5
      \qquad
      \lambda_{2}=2
  \]

  O autovalor $2$ possui multiplicidade algébrica dois
\end{resolution}

%------------------------------------------------------------------------------%
\begin{resolution}{Autoespaço para $\lambda=5$}
  Resolvemos
  \begin{align*}
      (B-5I)\vec{x}&=\vec{0}\\
      \begin{pmatrix}
          -2 & 1 & 1\\
          1 & -2 & 1\\
          1 & 1 & -2
      \end{pmatrix}
      \vetortri{x}{y}{z}
      &=\vetortri{0}{0}{0}
  \end{align*}

  O sistema equivalente fornece
  \begin{align*}
      -2x+y+z&=0\\
      x-2y+z&=0\\
      x+y-2z&=0
  \end{align*}

  Subtraindo as equações
  \begin{align*}
      x&=y\\
      y&=z
  \end{align*}
\end{resolution}

%------------------------------------------------------------------------------%
\begin{resolution}{Autoespaço para $\lambda=5$}
  Portanto
  \[
      E_{5}=
      \operatorname{span}
      \left\{
          \vetortri{1}{1}{1}
      \right\}
  \]

  Escolhemos
  \[
      \vec{v}_{1}=\vetortri{1}{1}{1}
  \]

  Sua norma é
  \[
      \norm{\vec{v}_{1}}=\sqrt{3}
  \]

  O autovetor unitário é
  \[
      \vec{u}_{1}=
      \frac{1}{\sqrt{3}}
      \vetortri{1}{1}{1}
  \]
\end{resolution}

%------------------------------------------------------------------------------%
\begin{resolution}{Autoespaço para $\lambda=2$}
  Resolvemos
  \begin{align*}
      (B-2I)\vec{x}&=\vec{0}\\
      \begin{pmatrix}
          1 & 1 & 1\\
          1 & 1 & 1\\
          1 & 1 & 1
      \end{pmatrix}
      \vetortri{x}{y}{z}
      &=\vetortri{0}{0}{0}
  \end{align*}

  O sistema equivalente é
  \[
      x+y+z=0
  \]

  Escolhendo $x=s$ e $y=t$
  \[
      z=-s-t
  \]
\end{resolution}

%------------------------------------------------------------------------------%
\begin{resolution}{Autoespaço para $\lambda=2$}
  Portanto
  \begin{align*}
      E_{2}
      &=
      \left\{
          \vetortri{s}{t}{-s-t}
          \mid s,t\in\mathbb{R}
      \right\}\\
      &=
      \operatorname{span}
      \left\{
          \vetortri{1}{0}{-1}
          ,
          \vetortri{0}{1}{-1}
      \right\}
  \end{align*}

  Temos
  \[
      \dim(E_{2})=2
  \]
\end{resolution}

%------------------------------------------------------------------------------%
\begin{resolution}{Ortogonalização}
  Escolhemos os vetores
  \[
      \vec{v}_{2}=\vetortri{1}{0}{-1}
      \qquad
      \vec{v}_{3}=\vetortri{0}{1}{-1}
  \]

  Calculamos
  \begin{align*}
      \vec{v}_{2}\cdot\vec{v}_{3}
      &=-1\\
      \norm{\vec{v}_{2}}^{2}
      &=2
  \end{align*}

  Aplicamos Gram-Schmidt
  \begin{align*}
      \vec{w}_{2}
      &=\vec{v}_{2}\\
      \vec{w}_{3}
      &=\vec{v}_{3}
      -\frac{\vec{v}_{3}\cdot\vec{w}_{2}}
      {\norm{\vec{w}_{2}}^{2}}\vec{w}_{2}\\
      &=\vec{v}_{3}+\frac{1}{2}\vec{v}_{2}
  \end{align*}
\end{resolution}

%------------------------------------------------------------------------------%
\begin{resolution}{Ortogonalização}
  Substituindo
  \begin{align*}
      \vec{w}_{3}
      &=
      \vetortri{0}{1}{-1}
      +\frac{1}{2}\vetortri{1}{0}{-1}\\
      &=
      \vetortri{\frac{1}{2}}{1}{-\frac{3}{2}}
  \end{align*}

  Podemos multiplicar por $2$ sem alterar a direção
  \[
      \vec{w}_{3}=\vetortri{1}{2}{-3}
  \]

  Verificamos
  \begin{align*}
      \vec{w}_{2}\cdot\vec{w}_{3}
      &=1-3=0
  \end{align*}
\end{resolution}

%------------------------------------------------------------------------------%
\begin{resolution}{Normalização}
  Calculamos as normas
  \begin{align*}
      \norm{\vec{w}_{2}}
      &=\sqrt{1^{2}+0^{2}+(-1)^{2}}\\
      &=\sqrt{2}\\
      \norm{\vec{w}_{3}}
      &=\sqrt{1^{2}+2^{2}+(-3)^{2}}\\
      &=\sqrt{14}
  \end{align*}

  Os vetores unitários são
  \begin{align*}
      \vec{u}_{2}
      &=\frac{1}{\sqrt{2}}
      \vetortri{1}{0}{-1}\\
      \vec{u}_{3}
      &=\frac{1}{\sqrt{14}}
      \vetortri{1}{2}{-3}
  \end{align*}
\end{resolution}

%------------------------------------------------------------------------------%
\begin{resolution}{Base Ortonormal}
  Obtivemos a base ortonormal de autovetores
  \begin{align*}
      \vec{u}_{1}
      &=\frac{1}{\sqrt{3}}
      \vetortri{1}{1}{1}\\
      \vec{u}_{2}
      &=\frac{1}{\sqrt{2}}
      \vetortri{1}{0}{-1}\\
      \vec{u}_{3}
      &=\frac{1}{\sqrt{14}}
      \vetortri{1}{2}{-3}
  \end{align*}

  O primeiro vetor pertence a $E_{5}$

  Os outros dois pertencem a $E_{2}$
\end{resolution}

%------------------------------------------------------------------------------%
\begin{resolution}{Matriz ortogonal}
  Organizamos os autovetores nas colunas de $P$
  \[
      P=
      \begin{pmatrix}
          \frac{1}{\sqrt{3}} &
          \frac{1}{\sqrt{2}} &
          \frac{1}{\sqrt{14}}\\[0.2cm]
          \frac{1}{\sqrt{3}} &
          0 &
          \frac{2}{\sqrt{14}}\\[0.2cm]
          \frac{1}{\sqrt{3}} &
          -\frac{1}{\sqrt{2}} &
          -\frac{3}{\sqrt{14}}
      \end{pmatrix}
  \]

  Como as colunas formam uma base ortonormal
  \[
      P^tP=I
  \]
\end{resolution}

%------------------------------------------------------------------------------%
\begin{resolution}{Matriz diagonal}
  Os autovalores correspondentes aos vetores das colunas de $P$ são
  \[
      5\qquad 2\qquad 2
  \]

  Portanto
  \[
      D=
      \begin{pmatrix}
          5 & 0 & 0\\
          0 & 2 & 0\\
          0 & 0 & 2
      \end{pmatrix}
  \]

  Assim
  \[
      B=PDP^t
  \]
\end{resolution}

%------------------------------------------------------------------------------%
\begin{resolution}{Resposta}
  A matriz ortogonal é
  \[
      P=
      \begin{pmatrix}
          \frac{1}{\sqrt{3}} &
          \frac{1}{\sqrt{2}} &
          \frac{1}{\sqrt{14}}\\
          \frac{1}{\sqrt{3}} &
          0 &
          \frac{2}{\sqrt{14}}\\
          \frac{1}{\sqrt{3}} &
          -\frac{1}{\sqrt{2}} &
          -\frac{3}{\sqrt{14}}
      \end{pmatrix}
  \]

  A matriz diagonal é
  \[
      D=
      \begin{pmatrix}
          5 & 0 & 0\\
          0 & 2 & 0\\
          0 & 0 & 2
      \end{pmatrix}
  \]

  Portanto
  \[
      B=PDP^t
  \]
\end{resolution}

%------------------------------------------------------------------------------%
